\appendix
\chapter{Annexes}

\section{Captures dashboard}
Les figures suivantes completent la documentation de l'interface Ariel \logminer.

\begin{figure}[H]
  \centering
  \includegraphics[width=\textwidth]{dashboard_longue_vue.png}
  \caption{Vue longue du dashboard.}
\end{figure}

\section{Tableaux complementaires}
Les tableaux Markdown complets sont fournis dans le dossier \texttt{tables/} du projet LaTeX et dans le pack de redaction final. Ils couvrent notamment les faux positifs, la comparaison Wazuh, les campagnes superviseur, la resilience, Redis Streams, MQTT et les ablations de routage.

\section{Commandes de reproduction}
Les commandes principales utilisees pour reproduire les resultats sont les suivantes. Elles sont presentees sur plusieurs lignes afin de rester lisibles dans le document imprime :

\begin{itemize}
  \item \path{python scripts/benchmark_realtime_workflow.py}
  \item \path{python scripts/run_resource_campaign.py}
  \item options principales : \path{--cycles 30 --interval-sec 2 --max-mb 5}
  \item \path{python scripts/run_family_routing_ablation.py}
  \item \path{python scripts/generate_memoire_assets.py}
  \item \path{python scripts/export_memoire_figures.py}
\end{itemize}

\section{Depot GitHub}
Le depot GitHub associe au projet est :

\begin{center}
\url{https://github.com/andybitati/Memoire/tree/version_3}
\end{center}

Il permet de retrouver le code source, les scripts de traitement, les ressources experimentales versionnees et les evolutions futures du prototype. Les articles scientifiques seront ajoutes plus tard lorsqu'ils seront rediges et finalises.

\section{Fichiers de preuves}
Les preuves experimentales principales sont conservees sous forme de CSV :
\begin{itemize}
  \item \texttt{realtime\_workflow\_benchmark\_20260604.csv} ;
  \item \texttt{resource\_campaign.csv} ;
  \item \texttt{family\_routing\_ablation.csv} ;
  \item \texttt{standard\_tools\_wazuh\_logminer\_overlap.csv} ;
  \item \texttt{random\_forest\_linux\_auth\_metrics.csv} ;
  \item \texttt{random\_forest\_network\_cicids\_metrics.csv} ;
  \item \texttt{random\_forest\_unsw\_80\_20\_metrics.csv}.
\end{itemize}

\section{Regles de prudence}
Pour toute lecture des resultats, les regles suivantes doivent etre conservees :
\begin{itemize}
  \item les anomalies non supervisees sont des anomalies candidates ;
  \item la baseline fail2ban-like n'est pas une execution officielle de fail2ban ;
  \item Wazuh/OSSEC restent des references operationnelles plus matures ;
  \item la distribution multi-machine stricte est une perspective ;
  \item l'auto-apprentissage continu est une extension future.
\end{itemize}

\section{Checklist de depot}
Avant le depot final, les elements suivants doivent etre verifies :

\begin{enumerate}
  \item orthographe officielle de tous les noms figurant sur la page de garde ;
  \item coherence entre le titre de la page de garde, le resume et le document directeur ;
  \item presence des figures et captures dans les dossiers \texttt{figures} et \texttt{captures} ;
  \item compilation Overleaf sans erreur bloquante ;
  \item coherence des numeros de figures et tableaux apres compilation ;
  \item verification des valeurs numeriques citees dans le chapitre des resultats ;
  \item presence de la bibliographie et des citations principales ;
  \item relecture finale des accents, noms propres et formulations administratives.
\end{enumerate}

\section{Fichiers de preuve recommandes}
Les fichiers de preuve a conserver avec le memoire sont les CSV de benchmark, les tableaux Markdown, les captures du dashboard, les figures exportees et les scripts de generation. Ils permettent de justifier les valeurs reprises dans le texte et de montrer que les resultats ne sont pas seulement descriptifs, mais rattaches a des artefacts reproductibles.

Les scripts les plus importants sont ceux qui generent les assets du memoire, executent les benchmarks de latence, lancent la campagne CPU/RAM, testent le routage par famille et consolident les metriques supervisees. Les fichiers sources du pipeline, de l'API, du routeur, du detecteur, du correlateur et de l'audit completent ces preuves.
